%! TeX root = ./main.tex
% vim: set spelllang=fr:
\chapter{Réconciliation bilatérale}

L’évolution de certaines unités, que nous appellerons \emph{dépendantes}, peut être liée à celle d’un hôte auquel elles sont attachées, formant une hiérarchie évolutive.
Pour étudier l’histoire de telles hiérarchies, l’arbre phylogénétique de l’unité dépendante peut être complété en déterminant, pour chacun de ses embranchements, s’il a pour origine un événement affectant l’unité dépendante ou l’hôte.

Les gènes sont un exemple d’unités dépendantes de l’évolution de leurs espèces hôtes.
Lorsqu’un gène se sépare en deux copies à la suite d’une duplication --- un événement l’affectant directement --- ces copies sont dites \emph{paralogues}.
Lorsque cette séparation est due à une spéciation de son hôte, les copies sont dites \emph{orthologues}~\cite{Fitch1970DistinguishingHomologousAnalogous}.
Savoir si deux gènes sont orthologues ou paralogues peut contribuer à déterminer leur fonction.
La fonction d’un gène a en effet tendance à être mieux conservée parmi les copies orthologues que parmi les copies paralogues~\cite{Lynch2001ProbabilityPreservationNewly}.
Par ailleurs, les gènes orthologues renseignent directement sur l’arbre phylogénétique des espèces hôtes.
Ainsi, l’arbre phylogénétique d’une famille de gènes tous orthologues donne une estimation de l’arbre des espèces hôtes.

La méthode de \emph{réconciliation}, proposée par \citeauthor{Goodman1979FittingGeneLineage} en~\citeyear{Goodman1979FittingGeneLineage}~\cite{Goodman1979FittingGeneLineage}, puis formalisée par~\citeauthor{Page1994MapsTreesCladistic} en~\citeyear{Page1994MapsTreesCladistic}~\cite{Page1994MapsTreesCladistic}, offre un cadre permettant d’inférer le statut de chaque divergence dans l’arbre phylogénétique d’une unité dépendante.
Ce chapitre est consacré à la présentation de cette méthode.

\section{Plongement d’arbres phylogénétiques}

Chaque embranchement de l’arbre phylogénétique d’une unité dépendante est dû soit à sa divergence directe, soit à la divergence de son hôte.
Dans une phylogénie à \(n\) embranchements, il peut ainsi exister jusqu’à \(2^n\) possibilités pour affecter chaque embranchement à son type.
La topologie de la phylogénie des hôtes peut être utilisée comme un guide pour évaluer ces affectations et sélectionner les plus plausibles.

Supposons que la phylogénie \(T\) d’une famille dépendante et celle de ses hôtes \(S\) soient connues.
Chaque feuille de \(T\), correspondant à une unité dépendante actuelle, est associé à une feuille de \(S\), son hôte.
Lorsque la divergence d’un hôte cause un embranchement \(v\) dans \(T\), ce dernier se rapporte à un embranchement correspondant dans~\(S\).
Au contraire, si \(v\) se produit indépendamment de l’évolution de l’hôte, donc entre deux embranchements de \(S\), il se rapporte à une branche de~\(S\).
Cette association de \(T\) avec \(S\) est appelée \emph{réconciliation}.
Pour être valide, une réconciliation ne doit pas \enquote{remonter dans le temps}, c’est-à-dire qu’un descendant de \(v\) ne doit pas être associé à un ancêtre de l’associé de \(v\)~(voir~\autoref{figure:invalid-rec}).

Formellement, une réconciliation est un plongement de \(T\) dans \(S\), au sens défini ci-après.
Pour simplifier, lorsque \(v\) doit s’associer à une branche de \(S\), il est associé à sa terminaison.

\begin{definition}[Réconciliation]\label{definition:rec}
    Soient \(T\) et \(S\) deux arbres phylogénétiques et \(\llmap : \tleavs(T) \onto \tleavs(S)\) une application qui associe chaque feuille de \(T\) à son hôte dans \(S\).
    Une \emph{réconciliation} de \(T\) avec \(S\) est une application \(\rec : \tnodes(T) \onto \tnodes(S)\) qui étend \(\llmap\).
    Une réconciliation est \emph{valide} si aucune paire \((v, w) \in \tnodes(T)^2\) avec \(v \tancs w\) n’est telle que \(\rec(v) \tdesc \rec(w)\).
\end{definition}

\begin{figure}[htb]
    \capstart
    \centering
    \begin{tikzpicture}
    %! TeX root = ../main.tex
% vim: set spelllang=fr:

\tikzset{
    % Reconciliation edges
    associate edge/.style={
        draw, thick, #1,
        shorten <=4pt, shorten >=3pt,
        ->, >=Latex,
    },
    overlaid associate edge/.style={
        draw, line width=4pt, white,
        shorten <=8pt, shorten >=16pt,
        postaction={associate edge=#1},
    },
    associate/.style={every edge/.style={associate edge=#1}},
    overlaid associate/.style={every edge/.style={overlaid associate edge=#1}},
    % Tree layout
    /graph drawing/evolutionary unit length=.66cm,
    /graph drawing/sibling distance=.4cm,
    % Predefined associate tree
    graphs/declare={ttree}{
        l3s1 / [coordinate] -- {
            l2s2 / [coordinate] -- {
                l1s4 / 3,
                l1s3 / 1,
            },
            l2s1 / [coordinate] -- {
                l1s2 / 2,
                l1s1 / 1,
            },
        },
    },
    % Predefined host tree
    graphs/declare={stree}{
        l3s1 / [coordinate] -- {
            l1s3 / 3 [> length=2],
            l2s1 / [coordinate] -- {
                l1s2 / 2,
                l1s1 / 1,
            },
        },
    },
}

\colorlet{reccolor1}{blue!60!white}
\colorlet{reccolor2}{red!80!white}
\colorlet{reccolor3}{green!60!black}


    \coordinate (l0) at (0, -.6);
    \coordinate (l1) at (0, 0);
    \coordinate (l2) at (0, .9);
    \coordinate (l3) at (0, 1.5);

    \makeTtree{t1}{0,0}
    \makeTtree{t2}{8,0}
    \makeStree{s}{4.5,0}

    \colorlet{s1}{blue!60!white}
    \colorlet{s2}{red!80!white}

    \path[associate=s1] (t1-l3-1) edge[bend left=10] node[above, midway] {\(\rec_1\)} (s-l3-1);
    \path[associate=s1] (t1-l2-1) edge[bend left=15] (s-l2-1);
    \path[overlaid associate=s1] (t1-l2-2) edge[out=0, in=180] (s-l3-1);

    \path[overlaid associate=s2] (t2-l3-1) edge[bend left=10] (s-l2-1);
    \path[overlaid associate=s2] (t2-l2-1) edge[bend left=10] (s-l2-1);
    \path[overlaid associate=s2] (t2-l2-2) edge[out=170, in=0] node[above, midway, s2] {\(\rec_2\)} (s-l3-1);
\end{tikzpicture}

    \caption{
        Deux réconciliations, \(\rec_1\) (valide) et \(\rec_2\) (invalide), entre deux phylogénies \(T\) et \(S\).
        Les associations entre embranchements sont représentées par des flèches.
        Pour alléger la figure, les associations entre feuilles sont implicites entre les feuilles portant le même nom.
    }
    \label{figure:invalid-rec}
\end{figure}

Dans certaines hiérarchies évolutives, les unités dépendantes sont capables de migrer de leur hôte vers un autre hôte coexistant mais séparé.
Par exemple, en particulier chez les procaryotes, les gènes peuvent subir un transfert horizontal d’une espèce à une autre.
Lorsqu’un changement d’hôte se produit le long d’une branche \((v, w) \in \tedges(T)\), l’espèce \(\rec(w)\) à laquelle se rapporte la terminaison de la branche est \emph{séparée} de l’espèce \(\rec(v)\) associée à l’origine.
Une réconciliation qui ne contient pas de tel changement d’hôte est dite \emph{stricte}, car alors l’évolution de l’unité dépendante suit strictement celle de son hôte.

\begin{definition}[Réconciliation stricte ou libre]\label{definition:rec-strict}
    Une réconciliation \(\rec\) de \(T\) avec \(S\) est dite \emph{stricte} si toutes les paires
    \((v, w) \in \tnodes(T)^2\) sont telles que si \(v \tancs w\) alors \(\rec(v) \tancs \rec(w)\).
    Sinon, \(\rec\) est dite \emph{libre}.
\end{definition}
\begin{remark}
    Toutes les réconciliations strictes sont valides.
\end{remark}

Dans la \autoref{figure:invalid-rec}, la réconciliation \(\rec_1\) est une réconciliation stricte tandis que \(\rec_2\) est une réconciliation libre mais invalide.
La réconciliation \(\rec_3\) de la \autoref{figure:free-rec} est libre et valide.

\begin{figure}[htb]
    \capstart
    \centering
    \begin{tikzpicture}
    %! TeX root = ../main.tex
% vim: set spelllang=fr:

\tikzset{
    % Reconciliation edges
    associate edge/.style={
        draw, thick, #1,
        shorten <=4pt, shorten >=3pt,
        ->, >=Latex,
    },
    overlaid associate edge/.style={
        draw, line width=4pt, white,
        shorten <=8pt, shorten >=16pt,
        postaction={associate edge=#1},
    },
    associate/.style={every edge/.style={associate edge=#1}},
    overlaid associate/.style={every edge/.style={overlaid associate edge=#1}},
    % Tree layout
    /graph drawing/evolutionary unit length=.66cm,
    /graph drawing/sibling distance=.4cm,
    % Predefined associate tree
    graphs/declare={ttree}{
        l3s1 / [coordinate] -- {
            l2s2 / [coordinate] -- {
                l1s4 / 3,
                l1s3 / 1,
            },
            l2s1 / [coordinate] -- {
                l1s2 / 2,
                l1s1 / 1,
            },
        },
    },
    % Predefined host tree
    graphs/declare={stree}{
        l3s1 / [coordinate] -- {
            l1s3 / 3 [> length=2],
            l2s1 / [coordinate] -- {
                l1s2 / 2,
                l1s1 / 1,
            },
        },
    },
}

\colorlet{reccolor1}{blue!60!white}
\colorlet{reccolor2}{red!80!white}
\colorlet{reccolor3}{green!60!black}


    \coordinate (l0) at (0, -.6);
    \coordinate (l1) at (0, 0);
    \coordinate (l2) at (0, .9);
    \coordinate (l3) at (0, 1.5);

    \makeTtree{t3}{0,0}
    \makeStree{s}{4.5,0}

    \colorlet{s3}{green!60!black}

    \path[associate=s3] (t3-l3-1) edge[bend left=10] node[above, midway] {\(\rec_3\)} (s-l2-1);
    \path[associate=s3] (t3-l2-1) edge[bend left=10] (s-l2-1);
    \path[associate=s3] (t3-l2-2) edge[bend right=5] (s-l2-1);
\end{tikzpicture}

    \caption{
        Une réconciliation libre et valide entre deux phylogénies \(T\) et \(S\).
        Un changement d’hôte a lieu entre l’embranchement qui précède la feuille \emph{3} et la feuille \emph{3} dans \(T\).
    }
    \label{figure:free-rec}
\end{figure}

Comme le suggère le lemme suivant, l’espace des réconciliations strictes est considérablement plus restreint que l’espace général des réconciliations.

\begin{lemma}
    Si \(\rec\) est une réconciliation stricte de \(T\) avec \(S\) alors, pour tout embranchement \(v \in \tinner(T)\),
    \(\rec(v) \tancs \lca\left\{ \llmap(l) \suchthat l \in \tleavs(\subtree{T}{v}) \right\}\).
\end{lemma}
\vspace{-1em}
\begin{proof}
    Soit \(\rec\) une réconciliation de \(T\) avec \(S\) et supposons qu’un embranchement \(v \in \tinner(T)\) soit tel que
    \(\rec(v) \nottancs \lca\left\{ \llmap(l) \suchthat l \in \tleavs(\subtree{T}{v}) \right\}\).
    Il existe donc \(w \in \tleavs(\subtree{T}{v})\) tel que \(\rec(v) \nottancs \llmap(w)\)
    et donc tel que \(\rec(v) \not\tancs \rec(w)\).
    Donc \(\rec\) n’est pas stricte.
\end{proof}

\section{Modèles d’évolution gène--espèce}

Pour mesurer la plausibilité d’une réconciliation, il est nécessaire de poser un modèle d’évolution qui permette de l’interpréter d’un point de vue biologique.
Un tel modèle spécifie l’ensemble des événements qui peuvent affecter, conjointement ou séparément, l’évolution des unités dépendantes et de leurs hôtes.
Sous l’hypothèse d’un certain modèle, deux catégories de méthodes coexistent pour évaluer la plausibilité d’une réconciliation.

Dans les méthodes de \emph{parcimonie}, chaque réconciliation est associée à une valeur reflétant le minimum d’événements nécessaires pour l’interpréter dans le modèle choisi.
L’objectif est alors de rechercher les réconciliations minimisant cette valeur, c’est-à-dire les plus parcimonieuses.
Ces réconciliations ne reflètent pas l’histoire complète des unités étudiées, puisque l’évolution emprunte rarement le chemin le plus court.
En revanche, la méthode de parcimonie a l’avantage de sélectionner des réconciliations permettant d’expliquer les données observées sans ajouter d’hypothèses supplémentaires.

Dans les méthodes \emph{probabilistes}, une probabilité est attribuée à chaque événement du modèle.
\tk{Préciser et sourcer}

Dans la suite nous présentons les principaux modèles d’évolution stricte gènes--espèces.

\paragraph{Modèle duplication--perte (DL)}

\cite{Goodman1979FittingGeneLineage,Page1994MapsTreesCladistic}
Dans le modèle DL, les espèces n’évoluent que par des spéciations, et les gènes peuvent également diverger par duplication à l’intérieur de leur espèce, ou encore être perdus.
Une réconciliation stricte \(\rec\) s’interprète sous le modèle~DL comme suit.
Pour un embranchement \(v \in \tinner(T)\):
\begin{itemize}
    \item \(v\) correspond à un événement de spéciation seulement lorsque \(\rec(v)\) est un embranchement de \(S\) et si \(\rec(v\lchild) \tdesc \rec(v)\lchild\) et \(\rec(v\rchild) \tdesc \rec(v)\rchild\) ou si \(\rec(v\lchild) \tdesc \rec(v)\rchild\) et \(\rec(v\rchild) \tdesc \rec(v)\lchild\);
    \item si \(v\) correspond à un événement de spéciation, alors chaque embranchement sur les chemins \(\rec(v) \tpath \rec(v\lchild)\) et \(\rec(v) \tpath \rec(v\rchild)\), excluant \(\rec(v)\), \(\rec(v\lchild)\) et \(\rec(v\rchild)\), implique un événement de spéciation suivi d’une perte;
    \item si \(v\) correspond à un événement de duplication, alors chaque embranchement sur les chemins \(\rec(v) \tpath \rec(v\lchild)\) et \(\rec(v) \tpath \rec(v\rchild)\), incluant \(\rec(v)\) mais excluant \(\rec(v\lchild)\) et \(\rec(v\rchild)\), implique un événement de perte.
\end{itemize}

\cite{Chauve2009NewPerspectivesGene}
\cite{Gorecki2006DlsTreesModel}

% figure : un arbre de gènes, un arbre d’espèces, réconciliations

% figure : un arbre de gènes, voisinage de l’arbre par branch swap, minimum de réconciliation


\paragraph{Modèle duplication--régulation (DR)}


\paragraph{Modèle duplication--perte--coalescence (DLC)}

Tri incomplet de lignées (ILS)
\cite[p. 71]{Page1994MapsTreesCladistic}
\cite{Rasmussen2012UnifiedModelingGene}
\cite{Wu2013MostParsimoniousReconciliation}

\paragraph{Modèle duplication--transfert--perte (DTL)}

\cite[pp. 68--70]{Page1994MapsTreesCladistic}
\cite{Hallett2001EfficientAlgorithmsLateral}
\cite{Hallett2004SimultaneousIdentificationDuplications}
\cite{Tofigh2011SimultaneousIdentificationDuplications}
\cite{Bansal2012EfficientAlgorithmsReconciliation}

\paragraph{Modèle duplication--transfert--perte généralisé (DTLx)}

\textbf{Espèces non-échantillonnées (DTLx)}
\cite{Weiner2021ImprovedDuplicationTransfer}

\section{Calcul de réconciliations}

\paragraph{Inférence de l’arbre d’espèces}
\cite[pp. 72--73]{Page1994MapsTreesCladistic}
\cite{Guigo1996ReconstructionAncientMolecular}

La hiérarchie évolutive entre gènes et espèces rend difficile l’utilisation de l’arbre de gène pour estimer l’arbre d’espèces. La réconciliation peut cependant être utilisée pour remédier à ce problème.

\cite{Ma2000GeneTreesSpecies}

\paragraph{Correction de l’arbre de gènes}

Phénomène de convergence évolutive peut brouiller l’inférence de l’arbre de gènes à partir de séquences nucléotidiques.

probabiliste : \cite{Arvestad2004GeneTreeReconstruction}


\tk{autres techniques/modèles non mentionnés, autres reviews}
